\section{Introduction}
Scientific meta-analysis aggregates study results to summarize the current state of knowledge, and it helps us understand the degree of consensus in a scientific literature. However, classical meta-analysis does not reveal how we arrived at this consensus, nor does it identify which studies were particularly influential at the moment that they were published.

Intuitively, a study may influence consensus in different ways.  It might reduce uncertainty by presenting robust evidence from a large-scale study using an established methodology. Or it might cast doubt on previous estimates by introducing a new methodology or by reporting surprising results.

In this paper, we posit that answering these questions requires a sequential approach to meta-analysis, one where we consider the order in which the studies appeared to track how each one affected our understanding. Because we are tracking \textit{beliefs}---in the sense that a meta-analysis estimate characterizes the collective beliefs of a scientific community---we need to account for our uncertainty and how it changes over time.

To these ends, we propose the \textit{sequential meta-analysis research trace} (SMART). SMART takes in a sequence of studies that target the same unknown quantity. With these inputs, it produces a sequential Bayesian meta-analysis of the estimate and quantifies the contribution of each study to the state of knowledge. SMART aligns with the end results of classical meta-analysis, but also provides insight into the path taken to reach the current consensus.

SMART begins with an assumed Bayesian meta-analysis model, such as the classical random-effects meta-analysis model \citep{Higgins_Thompson_Spiegelhalter_2009}. Under the chosen model, it considers each study sequentially, quantifying each one's contribution by measuring its effect on the posterior. Specifically, we use the Wasserstein distance \citep{Villani_2008} to measure the change, which captures movement in both the mean estimate and the uncertainty surrounding it. SMART thus gives a sequential picture of how a series of studies have led to the current state of knowledge, and the relative contribution of each to that picture.

When evaluating a sequence of scientific studies, the methodology of a study can be as important as its specific empirical results. 
 To this end, we extend the classical random effects model to one that also accounts for which studies used the same methodology, e.g., propensity score matching or randomized trials. Our \textit{labeled random effects model} assumes that each methodology comes with an unknown bias, with a belief that more reliable methodologies have smaller biases. When incorporated into SMART, this new model allows for the possibility that studies using a superior methodology can make a larger contribution to the consensus estimate. In particular, it both informs our estimate of the biases of \textit{other} methodologies, as well as providing new estimates of the target itself.

 To demonstrate SMART, consider \citep{Card_Krueger_94}, Card and Krueger's seminal 1994 study on the effect of a minimum wage on employment. This study is part of a larger research literature that seeks to estimate this effect, and it was particularly important in its use of difference-in-differences together with a border discontinuity design to avoid confounding by changes in employment prospects.  Specifically, we focus on the beginning of the New Minimum Wage Research literature \citep{Ehrenberg_92}, a research program aimed at leveraging natural experiments to provide more credible estimates of the effect of the minimum wage on employment. Figure
\ref{fig:MW_figure} (Left) shows the estimates of this literature, as characterized by \citep{Neumark_Shirley_22}, leading up to Card and Krueger's study.
Using the labeled random effects model, Figure \ref{fig:MW_figure} (Right) shows the result of SMART, this literature's research trace. This trace tracks the sequential research contribution -- measured in terms of our proposed learning metric -- for every study as it enters the literature. 
While classical meta-analysis accords Card and Krueger's study little weight from a retrospective vantage point, SMART suggests that their study was highly influential.

Our approach contributes to the literature on sequential meta-analysis~\citep{Higgins_Whitehead_Simmonds_2011,Van_Der_Tweel_Bollen_2010,Spence_Steinsaltz_Fanshawe_2016,Bollen_Uiterwaal_Van_Vught_Van_Der_Tweel_2006,Kulinskaya_21}.  This literature has focused on methods that guarantee valid sequential inference \citep{Higgins_Whitehead_Simmonds_2011} or try to assess whether enough evidence has been collected for the results to be conclusive~\citep{Spence_Steinsaltz_Fanshawe_2016,Bollen_Uiterwaal_Van_Vught_Van_Der_Tweel_2006}. Recent work \citep{Kulinskaya_21} has analyzed the properties of sequential meta-analysis in the presence of temporal trends. 
In contrast, we provide a tool that allows us to track scientific consensus as it evolves.
Further, we extend the well-established random-effects model for meta-analysis \citep{Higgins_Thompson_Spiegelhalter_2009,Higgins_Whitehead_Simmonds_2011} to incorporate information and form inferences about the methodologies' biases.
This is similar to the mixed linear model developed in \citep{Raudenbush_1985} which incorporates study characteristics to model the variation among studies' effect sizes. Our model is a continuation of this line of thought and allows the inclusion of a scientific communities' beliefs about various methodologies' biases.
Tools for the formal assessment of randomized \citep{cochrane_ch8} and non-randomized \citep{cochrane_ch25} studies' risk of bias have been developed and are routinely used by practitioners. Our labeled random-effects model adds to the literature about incorporating prior knowledge about studies (such as information about their methodology) into meta-analysis to improve evidence synthesis~\citep{Li_94,Sutton_2001,Doi_Thalib_2008}.

We illustrate SMART by analyzing datasets drawn from two ends of the spectrum of methodological innovation: one literature is methodologically homogeneous, while the other experiences a noteworthy methodological advance.
SMART measures the relative influence of each study and helps illuminate the trajectory of scientific progress in the two cases. When the trajectory contains a methodological advance, SMART calls attention to scientific contributions that a conventional meta-analysis might overlook.